% Markdown SHA-256: e38ddfcbc2481dd0fa76aab9ff40cd3c1db90a2bc956c29f20b957160e1f78fd
\section*{Abstract}\label{abstract}
\addcontentsline{toc}{section}{Abstract}

An input-compatible tool update can admit requests absent from its previous contract. We study a consent-relative abstraction of this change: projections onto selected field names, together with realizable enum-advertised mutation atoms. These sets describe declared requests, not the server's actual authority. A four-sign classifier reports widening, narrowing, neutral changes, or a need for review. For a specified flat object fragment, we give pen-and-paper arguments that widening has a constructive witness and that admitting only neutral and narrowing updates preserves the abstract baseline grant. The implementation passes a 2,048-pair single-field model and a new 114,244-pair two-field model using an independent JSON Schema validator. Sixteen authored fixtures and two selected, source-pinned release cases illustrate its behavior. On seven authored pairs and both release cases, mcp-diff 0.1.0 reports no breaking or warning changes while this policy escalates. These are differences between decision procedures, not measured security misses. The revised artifact corrects a mutation-witness defect, an unsatisfiable-schema definition, and an inaccurate release extraction.

\section{Motivation and scope}\label{motivation-and-scope}

MCP exposes named tools with input schemas and allows servers to notify clients that their tool list has changed. The protocol also requires server-side access controls; a schema is not itself an authorization mechanism \citep{mcp-tools}.

Consider a declared field named \emph{mode}, initially restricted to \emph{read}. Adding \emph{write} can preserve every previously valid input while changing what the model can request. Conversely, adding a required parser field such as \emph{format} can invalidate old inputs without adding a selected permission coordinate. Compatibility and the abstraction studied here can therefore order an update differently.

We implement an advisory schema comparison and an optional CLI gate. We do not implement a permission system, infer arbitrary server behavior, prove that a release is malicious, or establish that a schema-valid request succeeds. A field called \emph{format} might control deletion in the actual implementation; our default name predicate would not capture that meaning. An unchanged schema can accompany arbitrary behavior changes.

The contribution is a precise, limited model and a tested implementation of its update rules. Novelty relative to the full literature, ecosystem detection accuracy, and operational effectiveness are not established by this artifact.

\section{Abstract request grants}\label{abstract-request-grants}

\subsection{Supported object fragment}\label{supported-object-fragment}

Let \(\mathcal J\) be finite JSON values representable by the implementation and \(\mathcal O\) its JSON objects. Comparisons operate on object arguments. A schema \(S\) consists of a finite declared field set \(F_S\), property domains \(D_S(f)\), required fields \(R_S\subseteq F_S\), and an open/closed flag.

Each property domain is a JSON type or nonempty union of types, optionally intersected with a finite enum. Missing type denotes all represented JSON types; missing enum imposes no finite restriction. The empty enum denotes an empty domain. Integer values are included in the number type. Structured enum values use structural equality. JSON numbers outside the runtime's faithful numeric representation are not covered.

An object \(v\) is accepted exactly when:

\begin{enumerate}
\def\labelenumi{\arabic{enumi}.}
\tightlist
\item
  Every field in \(R_S\) occurs in \(v\).
\item
  Each declared field occurring in \(v\) has a value in \(D_S(f)\).
\item
  A closed schema has no undeclared fields in \(v\).
\end{enumerate}

Write this set as \(A(S)\). The schema is satisfiable iff every required property domain is nonempty; optional empty domains are permitted. Absent \emph{additionalProperties} means open. An absent root type is interpreted on the object-only argument domain, not as containment over arbitrary JSON.

These rules follow the relevant JSON Schema assertions \citep{json-validation, json-core}. We treat descriptions, titles, defaults, examples, and the recognized annotation keywords as non-asserting. A default does not fill in a missing value during this analysis. The implementation also ignores \emph{x-} extension keywords; the model assumes they carry no validation semantics.

Nested validation, combinators, references, patterns, numeric/string bounds, dependent constraints, schema-valued additional properties, and unsupported keywords route a changed contract to review. Object and array \emph{values} remain supported through unconstrained types or finite enums. Malformed schemas and prototype-sensitive property names are also reviewed. Identical serialized schemas are skipped even if outside the fragment; this is absence of observed change, not validation of those schemas.

\subsection{Fixed predicates and two components}\label{fixed-predicates-and-two-components}

Fix a field predicate \(\lambda(f)\) and value predicate \(\mu(e)\) for an approval epoch. The projection \(\pi_\lambda(v)\) retains exactly the fields whose names satisfy \(\lambda\). Define the projected request surface \[
P_\lambda(S)=\{\pi_\lambda(v):v\in A(S)\}.
\]

Let \(E_S(f)\) denote an explicitly present enum. Define \textbf{realizable advertised mutation atoms} by \[
M_\mu(S)=\{(f,e): e\in E_S(f),\ \mu(e),\
                 \exists v\in A(S)\text{ with }v[f]=e\}.
\] An atom requires an accepted complete request, not just a locally type-valid enum constant. Consequently an unsatisfiable schema has both components empty. An unconstrained string field contributes no enum-advertised atom merely because it could accept the string \emph{write}.

The abstract grant and its order are \[
G(S)=(P_\lambda(S),M_\mu(S)),\qquad
(P,M)\leq(P',M')\ \Longleftrightarrow\ P\subseteq P'\land M\subseteq M'.
\] This is a deliberately chosen abstraction. It is not a claim that approving an MCP connection automatically grants these sets in an operating system or authorization service.

For the closed object whose only field is insensitive \emph{format}, changing its enum from \emph{read} to \emph{read,write} leaves the projection set equal to \(\{\{\}\}\) but adds \((\texttt{format},\texttt{write})\) to \(M_\mu\). The separate component captures this policy choice without making all enum fields sensitive.

\subsection{Implemented predicate instance}\label{implemented-predicate-instance}

The implementation uses a case-insensitive substring regex with 25 alternatives: allow, allowlist, command, cmd, directory, dir, endpoint, exec, file, filename, filepath, host, mode, namespace, path, permission, role, root, scope, secret, shell, token, uri, url, workspace.

The mutation predicate accepts strings whose lowercase form is one of: append, create, delete, execute, modify, mutate, patch, post, put, remove, run, send, update, upload, write. Atoms retain the original value: \emph{write} and \emph{WRITE} are distinct values, although both satisfy the predicate. These predicates are fixed in source, not configurable through the current classifier API.

A fresh sensitive name always exists outside a finite declared field set for this regex. The witness builder starts at \(\texttt{permission\_delta\_secret}\) and appends \(\texttt{\_secret}\) until the name is absent from both endpoints. Generalizing the arguments to another predicate requires a corresponding fresh-name construction. Changing predicates invalidates an approval epoch.

\section{Classification and enforcement}\label{classification-and-enforcement}

The classifier examines contract names first. Added names require review; removed names are narrowing. A malformed surviving contract or unsupported changed schema requires review. For supported schemas, it handles satisfiability before individual edits: impossible-to-possible is widening, possible-to-impossible is narrowing, and impossible-to-impossible is neutral.

For two satisfiable endpoints, its rules are:

{\def\LTcaptype{none} % do not increment counter
\begin{longtable}[]{@{}
  >{\raggedright\arraybackslash}p{(\linewidth - 2\tabcolsep) * \real{0.5000}}
  >{\raggedright\arraybackslash}p{(\linewidth - 2\tabcolsep) * \real{0.5000}}@{}}
\toprule\noalign{}
\begin{minipage}[b]{\linewidth}\raggedright
Change
\end{minipage} & \begin{minipage}[b]{\linewidth}\raggedright
Result
\end{minipage} \\
\midrule\noalign{}
\endhead
\bottomrule\noalign{}
\endlastfoot
Closed object becomes open & Widening with a fresh sensitive-key witness \\
Open object becomes closed & Narrowing \\
Required sensitive field becomes optional & Widening with an omitted-field witness \\
Optional field becomes required & Narrowing \\
Add/remove a sensitive field or mutation-bearing coordinate & Review \\
Add required insensitive, non-mutating field & Narrowing \\
Add optional or remove insensitive, non-mutating field & Neutral \\
Enum keyword appears or disappears & Review \\
Add an effective sensitive or mutating enum value & Widening \\
Remove an effective sensitive or mutating enum value & Narrowing \\
Change only insensitive, pure enum values & Neutral \\
Type change makes an explicit mutation value valid & Widening \\
Type change removes projected or mutation values & Narrowing \\
Sensitive type gains values without a widening rule & Review \\
Other insensitive carrier changes & Neutral \\
\end{longtable}
}

Requiredness, enum, and type edits are inspected independently. A single release may receive multiple signs. Rotation can remove one value and add another; widening means \textbf{non-containment of the new grant}, not that the entire old grant is a strict subset of the new one. Narrowing labels include compatibility contractions that leave the abstract grant equal. Review means the procedure abstains or invokes policy; it is not a mathematical undecidability claim.

The CLI command below exits unsuccessfully if any entry is widening or review:

\begin{verbatim}
mcp-observatory diff baseline.json current.json \
  --fail-on-permission-delta review
\end{verbatim}

The \emph{widening} threshold permits review-class changes and therefore does not implement the preservation policy below. Without the flag, results are advisory. The CLI does not itself stop a server or collect renewed consent; the calling host or CI workflow must enforce its decision.

\section{Correctness arguments for the fragment}\label{correctness-arguments-for-the-fragment}

These are pen-and-paper arguments about the specified abstraction and rules. They are not a proof-assistant verification of the TypeScript program. They assume complete, accurately paired contract sets; represented finite JSON values; fixed predicates; and the stated annotation/object semantics.

\subsection{Product lemma}\label{product-lemma}

For a satisfiable flat schema, constraints on different declared coordinates are independent. A projected object is realizable iff its sensitive coordinates satisfy their domains and requiredness, its undeclared sensitive keys respect closedness, and every required insensitive coordinate can be filled from its nonempty domain. Optional insensitive coordinates may be omitted. Moreover \((f,e)\in M_\mu(S)\) iff the schema is satisfiable, the field explicitly enumerates \(e\), the value satisfies its type, and \(\mu(e)\). Fill all other required fields independently to establish the reverse implication. For an unsatisfiable schema both sets are empty.

\subsection{Proposition 1: constructive widening}\label{proposition-1-constructive-widening}

Every widening entry for supported endpoints has an object accepted by the head and rejected by the base, demonstrating a new projection or mutation atom.

\emph{Impossible-to-possible:} any head witness projects to an element of a nonempty set; the base projection set is empty.

\emph{Closed-to-open:} add a fresh sensitive field to a valid head request. Its projection cannot be realized by the closed base.

\emph{Required-to-optional sensitive field:} omit that field and fill the head's remaining required fields. No base request can realize the resulting projection, which lacks a base-required sensitive coordinate.

\emph{Enum or carrier widening:} select a newly valid sensitive value, or an explicit newly valid mutation value when the field is insensitive. In the former case its projected coordinate is impossible in the base. In the latter case its mutation atom is absent from the base. Fill the other head-required fields using the product lemma.

These arguments use the actual endpoints. Simultaneous contractions on other coordinates do not invalidate a witness built from the head. A newly unsatisfiable head is handled earlier and receives no widening entry.

\subsection{Proposition 2: admitted updates do not expand the grant}\label{proposition-2-admitted-updates-do-not-expand-the-grant}

If the classifier emits no widening or review entry for supported endpoints, then \(G(S_{\rm head})\leq G(S_{\rm base})\).

An unsatisfiable head has the empty grant. An unsatisfiable base with a satisfiable head would have emitted widening. Otherwise both endpoints are satisfiable. The head cannot open a closed base, add/remove a sensitive coordinate, make a required sensitive field optional, or gain an effective sensitive value: each such edit emits widening or review. Thus every head projection can be completed to a base request, filling any removed or changed insensitive required coordinates from the base's nonempty domains.

For mutation atoms, enum presence cannot change without review. No newly declared mutation coordinate can pass review, and no existing coordinate can gain a type-valid mutating enum value without widening. Every realized head atom therefore occurs in the base, by the product lemma. The two inclusions establish the result directly, without assuming an ordering of intermediate atomic edits that might become unsatisfiable.

\subsection{Proposition 3: neutral invariance}\label{proposition-3-neutral-invariance}

If the endpoints are supported and every emitted entry is neutral, or there are no entries, their grants are equal. Both impossible endpoints give empty grants. For satisfiable endpoints the previous argument gives one inclusion. With no narrowing entries, neither closedness nor sensitive requiredness can strengthen, and no effective sensitive value or mutation atom can disappear. Edits confined to insensitive pure coordinates preserve realizability by the product lemma. This gives the reverse inclusion. Syntactically identical inputs have identical abstractions under the fixed semantics.

\subsection{Corollary: preservation within an approval epoch}\label{corollary-preservation-within-an-approval-epoch}

Start with an approved abstract grant \(G_0\). If each transition blocks widening and review results and otherwise admits the head, Proposition 2 and transitivity give \(G_t\leq G_0\) for every admitted version. For a contract set, index both components by contract name; removal contracts the set, while additions require review. Renewed approval starts a \textbf{new} epoch with a new baseline; it does not preserve containment in the old epoch automatically.

This corollary requires enforcement of every transition and trustworthy contract identity/collection. It proves no corresponding inclusion between the servers' actual side effects, credentials, or resource access.

\subsection{Implementation cost}\label{implementation-cost}

The old draft's total linear-time claim is withdrawn. If there are \(W\) widening records and \(R\) required fields, materializing a complete request for each witness can output \(\Theta(WR)\) field occurrences, even with primitive enums. Structured enum canonicalization also requires recursion and key sorting. The implementation performs local domain comparisons rather than general schema containment, but we report neither a linear bound for total output nor a controlled performance benchmark.

\section{Validation and descriptive comparison}\label{validation-and-descriptive-comparison}

\subsection{Regression and finite-model checks}\label{regression-and-finite-model-checks}

The retained one-field model checks 2,048 pairs: two field labels, endpoint requiredness and closedness, two type carriers, and four enum configurations.

The added two-field model independently enumerates accepted requests using AJV 8.20.0. Each of \emph{path} and \emph{format} is absent, or optional/required with one of six domains: unconstrained; empty string enum; \emph{read}; \emph{read,png,write}; number with enum \emph{0,write}; or string/number with enum \emph{0,write}. Together with open/closed objects this gives 338 schemas and \(338^2=114{,}244\) ordered endpoint pairs.

The request grid contains absence, the empty string, \emph{read}, \emph{png}, \emph{write}, 0, 0.5, false, null, an empty array, and an empty object for each coordinate, plus an absent/present fresh sensitive key. AJV rejects an empty enum during compilation, so the oracle replaces only that property schema with \emph{false}, which has the same empty-domain semantics. It uses no classifier domain helpers.

For every admitted pair it checks componentwise non-expansion, and for every neutral-only pair it checks equality. Every widening witness is independently checked for head acceptance, base rejection, and a new projection or realized mutation atom. All pairs pass. The grid is finite: it is neither exhaustive over all JSON nor a mechanized proof of the propositions.

This revision exposed and corrected two issues. The earlier mutation definition counted optional \emph{write} even when another required field made every request impossible. Requiring whole-request realizability repairs that definition. Separately, an insensitive enum expanding from \emph{read} to \emph{read,png,write} previously returned a \emph{png} witness; that showed new acceptance but no new abstract atom. The implementation now selects \emph{write}. These are distinct changes to the model and implementation.

\subsection{Authored conformance fixtures}\label{authored-conformance-fixtures}

The 16 fixtures and expectations were authored with the rules. They all match their expected sign presence: seven seeded scenarios, eight synthetic coverage cases, and one enum rotation. They produce four widening, five review, four narrowing, and seven neutral records; some pairs have multiple records. They do not cover every possible rule or interaction.

{\def\LTcaptype{none} % do not increment counter
\begin{longtable}[]{@{}lr@{}}
\toprule\noalign{}
Highest-priority action & Authored pairs \\
\midrule\noalign{}
\endhead
\bottomrule\noalign{}
\endlastfoot
Widening: would block & 4 \\
Review: would queue & 5 \\
Neutral/narrowing: pass & 7 \\
No records & 0 \\
\end{longtable}
}

All 16 inputs change a contract. A binary-change policy would flag 16; the review threshold flags 9 (43.75\% fewer), and the widening-only threshold flags 4 (75\% fewer). These are routing counts for a constructed mixture, not false-positive reductions. Widening-only also admits cases for which the preservation argument supplies no guarantee.

\subsection{Selected release cases}\label{selected-release-cases}

Six upstream files were fetched at immutable commits and checked by SHA-256. The extraction script runs the reviewed, hash-pinned Notion converter and reads GitHub's released snapshots. Schema annotations and unused definitions are removed; remaining converter/snapshot validation keywords are preserved. The artifact contains only selected changed contracts, not complete release inventories or captures from live servers.

\textbf{Notion v2.1.0 to v2.3.1.} The earlier OpenAPI document has no page-Markdown path. The head converter generates retrieve and update operations, and its proxy advertises them as \emph{API-retrieve-page-markdown} and \emph{API-update-page-markdown}. Both additions produce review records. The current extraction preserves the converter's format, defaults, nested schemas, and anyOf string fallbacks; the old hand-written extraction did not reproduce these details. The new-tool rule does not claim containment for those nested schemas \citep{notion-parser, notion-proxy}.

\textbf{GitHub v0.7.0 to v0.8.0.} The released get-file-contents snapshots change required fields from \emph{owner,repo,path} to \emph{owner,repo}, adding the default slash to \emph{path}. The associated source change confirms the intended root default \citep{github-path}. The classifier emits widening, with \emph{owner} and \emph{repo} present and \emph{path} omitted. This is a schema acceptance witness; the empty placeholder repository values need not correspond to a successful API call. The case establishes a change in omission semantics, not a vulnerability or a newly accessible resource.

\subsection{Pinned mcp-diff comparison}\label{pinned-mcp-diff-comparison}

We ran the installed mcp-diff 0.1.0 classification engine, whose source hash matches its published wheel \citep{mcp-diff}. Each subprocess checks both version and module hash. The runner supplies empty tool descriptions to both sides and compares the same normalized input schemas. It recomputes this classifier's actions directly; it does not trust a stale result file.

{\def\LTcaptype{none} % do not increment counter
\begin{longtable}[]{@{}lll@{}}
\toprule\noalign{}
Pair & mcp-diff result & This policy \\
\midrule\noalign{}
\endhead
\bottomrule\noalign{}
\endlastfoot
rugpull-01: mode enum grows & Compatible & Widening \\
rugpull-02: closed object opens & Compatible & Widening \\
rugpull-03: required scope relaxed & Compatible & Widening \\
rugpull-04: token field added & Compatible & Review \\
multi-01: tool added & Compatible & Review \\
limit-01: enum rotation & Compatible & Widening \\
combo-02: multi-tool addition & Compatible & Review \\
Selected Notion additions & Compatible & Review \\
Selected GitHub path change & Compatible & Widening \\
\end{longtable}
}

Here \emph{compatible} means this pinned engine emitted neither breaking nor warning records. It is not proof of acceptance-set containment. In particular, version 0.1.0 does not inspect enum values, additional-properties changes, or requiredness changes on existing parameters. Its compatible result for the enum rotation therefore also reflects implementation coverage, not merely opposite mathematical orderings. New tools produce informational records, not an identical result.

One authored pair goes the other way: adding required insensitive \emph{format} is breaking under mcp-diff and admitted with narrowing under this policy. The source supports these version-specific observations; we make no claim about newer versions or comparative security quality. All disagreements are stored as descriptive outcomes, without the former ``consent miss'' label.

\section{Reproducibility and limitations}\label{reproducibility-and-limitations}

The source, commands, frozen inputs, hashes, baseline wheel pin, and extraction procedure are documented in the corpus extraction runbook. The Markdown is the canonical manuscript text; the LaTeX body is generated from it, and named/anonymous-author PDFs share that body. The anonymous-author variant is not certified to satisfy a particular venue's blind-review rules.

The theoretical guarantee is relative to the declared abstraction and trusted update handling. It excludes schema-independent behavior, semantic meanings not captured by the predicates, invalid collection/identity matching, arbitrary JSON Schema, and unfaithfully represented numbers. Only a finite implementation model was exhaustively checked. The two selected releases are illustrative, and the authored suite has no independent labels. No detection accuracy, ecosystem prevalence, user consent behavior, latency distribution, or ablation effect has been measured.

A useful next research step is a preregistered, independently sampled release corpus with separate labels for actual authorization effects and schema-level request changes. Until then, the supported result is a reproducible abstract change classifier with explicit review boundaries.
